\documentclass[12pt, a4paper]{report}

% ============================================================
% PACKAGES
% ============================================================
\usepackage[vietnamese]{babel}
\usepackage[utf8]{inputenc}
\usepackage[T1]{fontenc}
\usepackage{lmodern}
\usepackage{geometry}
\usepackage{graphicx}
\usepackage{booktabs}
\usepackage{tabularx}
\usepackage{array}
\usepackage{longtable}
\usepackage{multirow}
\usepackage{fancyhdr}
\usepackage{titlesec}
\usepackage{hyperref}
\usepackage{xcolor}
\usepackage{colortbl}
\usepackage{enumitem}
\usepackage{setspace}
\usepackage{caption}
\usepackage{subcaption}
\usepackage{float}
\usepackage{amsmath}
\usepackage{amssymb}
\usepackage{tikz}
\usepackage{pgfplots}
\pgfplotsset{compat=1.18}
\usetikzlibrary{
  shapes.geometric, shapes.symbols,
  arrows.meta, positioning, fit,
  backgrounds, shadows, decorations.pathreplacing,
  patterns, calc, matrix, chains
}
\usepackage{tcolorbox}
\tcbuselibrary{skins, breakable, theorems}
\usepackage{mdframed}
\usepackage{pifont}
\usepackage{fontawesome5}

% ============================================================
% PAGE GEOMETRY
% ============================================================
\geometry{
    top=2.5cm, bottom=2.5cm,
    left=3.0cm, right=2.0cm
}

% ============================================================
% COLORS
% ============================================================
\definecolor{ttuorange}{RGB}{211,84,0}
\definecolor{ttugold}{RGB}{243,156,18}
\definecolor{ttugreen}{RGB}{39,174,96}
\definecolor{ttublue}{RGB}{41,128,185}
\definecolor{ttured}{RGB}{192,57,43}
\definecolor{ttugray}{RGB}{80,80,80}
\definecolor{rowgray}{RGB}{245,245,245}
\definecolor{lightblue}{RGB}{214,234,248}
\definecolor{lightorange}{RGB}{253,237,226}
\definecolor{lightgreen}{RGB}{212,239,223}
\definecolor{lightred}{RGB}{250,219,216}
\definecolor{lightyellow}{RGB}{254,249,231}

% ============================================================
% TCOLORBOX STYLES
% ============================================================
\newtcolorbox{casebox}[1]{
  colback=lightorange, colframe=ttuorange,
  fonttitle=\bfseries\color{white}, coltitle=white, colbacktitle=ttuorange,
  title={#1}, boxrule=0.8pt, arc=4pt, left=6pt, right=6pt, top=4pt, bottom=4pt
}
\newtcolorbox{keybox}[1]{
  colback=lightblue, colframe=ttublue,
  fonttitle=\bfseries, coltitle=white, colbacktitle=ttublue,
  title={#1}, boxrule=0.8pt, arc=4pt
}
\newtcolorbox{warnbox}[1]{
  colback=lightred, colframe=ttured,
  fonttitle=\bfseries, coltitle=white, colbacktitle=ttured,
  title={#1}, boxrule=0.8pt, arc=4pt
}
\newtcolorbox{greenbox}[1]{
  colback=lightgreen, colframe=ttugreen,
  fonttitle=\bfseries, coltitle=white, colbacktitle=ttugreen,
  title={#1}, boxrule=0.8pt, arc=4pt
}

% ============================================================
% CHAPTER / SECTION FORMATTING
% ============================================================
\titleformat{\chapter}[block]
  {\normalfont\Large\bfseries\color{ttuorange}}
  {CHƯƠNG \thechapter:}{0.8em}{}
  [\vspace{-0.4em}\color{ttuorange}\rule{\linewidth}{1.5pt}\vspace{0.2em}]

\titlespacing*{\chapter}{0pt}{8pt}{12pt}

\titleformat{\section}
  {\normalfont\large\bfseries\color{ttuorange}}
  {\thesection}{0.8em}{}

\titleformat{\subsection}
  {\normalfont\normalsize\bfseries\color{ttugray}}
  {\thesubsection}{0.8em}{}

% ============================================================
% HEADER / FOOTER
% ============================================================
\pagestyle{fancy}
\fancyhf{}
\fancyhead[L]{\small\textit{Project Agent -- Trợ lý phân tích thị trường}}
\fancyhead[R]{\small\textit{Đề tài 4: RAG-Fusion}}
\fancyfoot[C]{\thepage}
\renewcommand{\headrulewidth}{0.4pt}

\fancypagestyle{plain}{
  \fancyhf{}
  \fancyhead[L]{\small\textit{Project Agent -- Trợ lý phân tích thị trường}}
  \fancyhead[R]{\small\textit{Đề tài 4: RAG-Fusion}}
  \fancyfoot[C]{\thepage}
  \renewcommand{\headrulewidth}{0.4pt}
}

% ============================================================
% MISC
% ============================================================
\setstretch{1.35}
\hypersetup{
  colorlinks=true, linkcolor=ttuorange,
  citecolor=ttuorange, urlcolor=ttublue,
  pdftitle={De tai 4 - Tro ly phan tich thi truong RAG-Fusion}
}
\setlist[itemize]{leftmargin=*, topsep=3pt, itemsep=2pt}
\setlist[enumerate]{leftmargin=*, topsep=3pt, itemsep=2pt}

% ============================================================
% TIKZ SHORTCUTS
% ============================================================
\tikzset{
  boxt/.style={rectangle, rounded corners=4pt, draw=#1, fill=#1!15,
    text width=3.2cm, align=center, minimum height=1cm, font=\small\bfseries},
  arrowt/.style={->, >=Stealth, thick, color=#1},
  process/.style={rectangle, rounded corners=3pt, draw=ttublue,
    fill=lightblue, minimum height=0.8cm, text width=3cm,
    align=center, font=\small},
  decision/.style={diamond, aspect=2, draw=ttuorange,
    fill=lightorange, minimum height=0.8cm, text width=2.5cm,
    align=center, font=\scriptsize},
  startstop/.style={rectangle, rounded corners=15pt, draw=ttugreen,
    fill=lightgreen, minimum height=0.8cm, text width=2.5cm,
    align=center, font=\small\bfseries},
  gnode/.style={rectangle, rounded corners=3pt, draw=#1, fill=#1!18,
    minimum height=0.9cm, text width=3.0cm, align=center, font=\footnotesize}
}

\begin{document}

% ============================================================
% TITLE PAGE
% ============================================================
\begin{titlepage}
\centering

\begin{tabular}{ccc}
  \includegraphics[width=3cm]{TTUlogo.png}
  & \hspace{2.5cm} &
  \includegraphics[width=3cm]{TTUITlogo.png}
\end{tabular}

\vspace{0.4cm}
{\Large\textbf{TRƯỜNG ĐẠI HỌC TÂN TẠO}} \\[3pt]
{\normalsize KHOA CÔNG NGHỆ THÔNG TIN}\\[3pt]
{\small Học phần: \textbf{PROJECT AGENT}}

\vspace{0.3cm}
{\color{ttuorange}\rule{\linewidth}{2pt}}
\vspace{1.2cm}

{\LARGE\textbf{\color{ttuorange}TIỂU LUẬN CUỐI KỲ}}\\[10pt]
{\Large\textbf{Đề tài 4}}\\[8pt]
{\large\textbf{TRỢ LÝ PHÂN TÍCH THỊ TRƯỜNG -- RAG-FUSION}}\\[6pt]
{\normalsize\textit{Market Analysis Assistant -- Reciprocal Rank Fusion}}

\vspace{1.5cm}

\begin{tabular}{ll}
  \textbf{Nhóm thực hiện:} & Nhóm F \\[5pt]
  \textbf{Thành viên:}
      & Đào Minh Hùng \quad -- \quad MSHV: 2632006 \\
      & Vương Thảo Nguyên -- MSHV: 2532013 \\
      & Phạm Thị Kim Ngọc -- MSHV: 2532005 \\
      & Nguyễn Khánh Hòa -- MSHV: 2532018 \\[5pt]
  \textbf{Giảng viên hướng dẫn:} & Trần Quốc Huy \\[5pt]
\end{tabular}

\vfill
{\color{ttuorange}\rule{\linewidth}{2pt}}\\[5pt]
{\normalsize Tây Ninh, tháng 9 năm 2026}
\end{titlepage}

% ============================================================
% MỤC LỤC
% ============================================================
\tableofcontents
\newpage

% ============================================================
% TÓM TẮT
% ============================================================
\chapter*{Tóm tắt}
\addcontentsline{toc}{chapter}{Tóm tắt}

Tiểu luận này trình bày \textbf{Trợ lý phân tích thị trường chứng khoán Việt Nam} xây dựng trên \textbf{LangGraph}, với kỹ thuật truy hồi trọng tâm là \textbf{RAG-Fusion}: phân rã một câu hỏi thành nhiều truy vấn con, truy hồi song song từng truy vấn, rồi hợp nhất thứ hạng bằng \textbf{Reciprocal Rank Fusion (RRF)}. Trọng tâm báo cáo là \textbf{luồng graph trạng thái} -- cách một câu hỏi ngôn ngữ tự nhiên chảy qua graph để thành báo cáo Markdown.

\medskip\noindent Các đóng góp chính:
\begin{enumerate}
  \item \textbf{Kiến trúc 3 tầng}: \texttt{conversation\_router} (phân loại, LLM-on-top) $\rightarrow$ graph LangGraph $\rightarrow$ intent modules + tools lấy dữ liệu thực.
  \item \textbf{Graph trạng thái đầy đủ}: \texttt{classify} $\rightarrow$ \texttt{clarify} $\rightarrow$ \texttt{decompose/fast-path} $\rightarrow$ \texttt{run\_subqueries} $\rightarrow$ \texttt{synthesize} $\rightarrow$ \texttt{critique} $\rightarrow$ \texttt{finalize}, với hai vòng lặp có trần (replan, critique).
  \item \textbf{RAG-Fusion}: graph 5 node \texttt{decompose} $\rightarrow$ \texttt{multi\_retrieve} $\rightarrow$ \texttt{rrf\_fuse} $\rightarrow$ \texttt{analyze} $\rightarrow$ \texttt{report}, hợp nhất BM25 + vector bằng RRF ($k=60$).
  \item \textbf{Ràng buộc trung lập}: cờ \texttt{STRICT\_NEUTRAL} cấm khuyến nghị mua/bán/nắm giữ trong bước tổng hợp.
\end{enumerate}

\medskip\noindent\textbf{Từ khóa:} RAG-Fusion, Reciprocal Rank Fusion, LangGraph, agent, BM25, Qdrant, Postgres, DeepSeek, RAGAS.

\newpage

% ============================================================
% CHƯƠNG 1: GIỚI THIỆU
% ============================================================
\chapter{Giới thiệu}

\section{Bối cảnh và Động lực}

Phân tích thị trường đòi hỏi trả lời những câu hỏi \textbf{phức hợp, đa góc nhìn}: diễn biến giá, tin tức ảnh hưởng, số liệu tài chính, so sánh nhiều mã. Một câu hỏi như \textit{``Đánh giá HPG hôm nay''} thực chất gồm nhiều intent con (giá, kỹ thuật, định giá, tin tức). RAG truyền thống -- một query, một lần truy hồi -- không đủ bao phủ một câu hỏi phức hợp.

\begin{casebox}{Vấn đề của truy hồi đơn lẻ}
Một query \textit{``HPG quý gần nhất''} không thể tối ưu đồng thời cho ``xu hướng giá'', ``tin ngành thép'', ``kết quả kinh doanh'' và ``vĩ mô''. Kết quả truy hồi lệch về một khía cạnh, bỏ sót phần còn lại -- độ phủ (recall) thấp.
\end{casebox}

Hệ thống giải quyết bằng hai cơ chế bổ trợ: (1) \textbf{decompose} -- tách câu hỏi thành các sub-task có intent riêng; (2) \textbf{RAG-Fusion} -- sinh nhiều truy vấn con và hợp nhất thứ hạng bằng RRF.

\section{Phạm vi và Mục tiêu}

\noindent\textbf{Đối tượng}: Trợ lý phân tích thị trường chứng khoán Việt Nam (cổ phiếu HOSE/HNX/UPCOM), triển khai theo \textbf{Đề tài 4 -- RAG-Fusion}.

\medskip\noindent\textbf{Mục tiêu cụ thể}:
\begin{enumerate}
  \item Dựng graph trạng thái LangGraph điều phối toàn bộ một lượt (turn): phân loại, làm rõ, phân rã, thu thập, tổng hợp, tự kiểm.
  \item Cài kỹ thuật RAG-Fusion: nhiều truy vấn con + RRF hợp nhất thứ hạng.
  \item Lấy dữ liệu thực qua tầng tools (giá, tin, định giá, vĩ mô, RAG) và tổng hợp trung lập.
  \item Chặn mọi vòng lặp (ngân sách turn) để chi phí và độ trễ bị giới hạn.
\end{enumerate}

\section{Cấu trúc Báo cáo}

\begin{figure}[H]
\centering
\begin{tikzpicture}[node distance=0.6cm, every node/.style={font=\small}]
  \node[startstop] (ch1) {Chương 1\\Giới thiệu};
  \node[process, right=0.5cm of ch1] (ch2) {Chương 2\\Kiến trúc\\tổng quan};
  \node[process, right=0.5cm of ch2] (ch3) {Chương 3\\Graph LangGraph\\(luồng chính)};
  \node[process, right=0.5cm of ch3] (ch4) {Chương 4\\RAG-Fusion \&\\RRF};
  \node[process, below=0.6cm of ch4] (ch5) {Chương 5\\Tầng tools \&\\dữ liệu};
  \node[process, left=0.5cm of ch5] (ch6) {Chương 6\\Cơ chế\\vận hành};
  \node[startstop, left=0.5cm of ch6] (ch7) {Chương 7\\Kết luận};

  \draw[arrowt=ttuorange] (ch1) -- (ch2);
  \draw[arrowt=ttuorange] (ch2) -- (ch3);
  \draw[arrowt=ttuorange] (ch3) -- (ch4);
  \draw[arrowt=ttuorange] (ch4) -- (ch5);
  \draw[arrowt=ttuorange] (ch5) -- (ch6);
  \draw[arrowt=ttuorange] (ch6) -- (ch7);
\end{tikzpicture}
\caption{Cấu trúc tổng thể của báo cáo}
\end{figure}

\newpage

% ============================================================
% CHƯƠNG 2: KIẾN TRÚC TỔNG QUAN
% ============================================================
\chapter{Kiến trúc tổng quan}

\section{Ba tầng xử lý}

Một lượt đi qua ba tầng, trong đó \textbf{LangGraph là tầng giữa} -- nơi toàn bộ trạng thái được điều phối (chi tiết Chương 3):

\begin{enumerate}
  \item \textbf{Tầng 1 -- \texttt{conversation\_router}}: một lần gọi LLM ép chọn tool, đổ kết quả (\texttt{intent}, \texttt{ticker}) vào \texttt{AgentState}.
  \item \textbf{Tầng 2 -- graph LangGraph}: điều phối toàn bộ luồng \texttt{classify} $\rightarrow$ \texttt{clarify} $\rightarrow$ \texttt{decompose} $\rightarrow$ \texttt{run\_subqueries} $\rightarrow$ \texttt{synthesize} $\rightarrow$ \texttt{critique} $\rightarrow$ \texttt{finalize}.
  \item \textbf{Tầng 3 -- intent modules + tools}: \texttt{gather\_data}/\texttt{run} lấy dữ liệu thực (giá, chỉ báo, định giá, tin, vĩ mô, RAG).
\end{enumerate}

\section{Stack công nghệ}

\begin{table}[H]
\centering
\caption{Stack công nghệ}
\small
\begin{tabularx}{\textwidth}{lX}
\toprule
\textbf{Layer} & \textbf{Công nghệ} \\
\midrule
LLM & DeepSeek (default) $\cdot$ Anthropic $\cdot$ OpenAI $\cdot$ Ollama $\cdot$ Gemini -- qua \texttt{create\_client()} \\
\rowcolor{rowgray}
Embedding & bge-m3 (dim 1024) $\cdot$ nomic-embed-text (dim 768) \\
Vector store & Qdrant \\
\rowcolor{rowgray}
Relational DB & Postgres (\texttt{financial\_facts}, \texttt{stock\_prices}, \texttt{news\_articles}, \ldots) \\
Object storage & MinIO (PDF BCTC) \\
\rowcolor{rowgray}
Orchestration & Dagster (schedules + sensor) \\
Graph & LangGraph \\
\rowcolor{rowgray}
Reranker & CrossEncoder ms-marco-MiniLM-L-6-v2 \\
Observability & Langfuse (tracing) \\
\bottomrule
\end{tabularx}
\end{table}

\begin{keybox}{LLM Factory -- create\_client()}
Mọi node gọi model qua \texttt{from llm.factory import create\_client} -- không import provider trực tiếp. \texttt{create\_client()} đọc \texttt{LLM\_PROVIDER} từ \texttt{.env} (default \texttt{deepseek}) và trả đúng \texttt{LLMClient}. Đổi provider = một dòng \texttt{.env}, không đổi code.
\end{keybox}

\section{Nguồn dữ liệu}

\begin{table}[H]
\centering
\caption{Nguồn dữ liệu và nơi lưu trữ}
\small
\begin{tabularx}{\textwidth}{lXXl}
\toprule
\textbf{Nguồn} & \textbf{Nội dung} & \textbf{Lưu ở đâu} & \textbf{Tag} \\
\midrule
PDF BCTC (HPG) & BCTC kiểm toán & Qdrant \texttt{hpg\_b7\_structural\_meta} & RAG corpus \\
\rowcolor{rowgray}
CafeF / VnExpress RSS & Tin tài chính, thép, TTCK & Postgres + Qdrant \texttt{news\_chunks} & TIN TỨC \\
vnstock Finance & KQKD, CDKT, LCTT & Postgres \texttt{financial\_facts} & GIÁ LỊCH SỬ \\
\rowcolor{rowgray}
vnstock Prices & Giá đóng cửa điều chỉnh & Postgres \texttt{stock\_prices} & GIÁ LỊCH SỬ \\
Tavily web search & Kết quả real-time & Không lưu & WEB \\
\rowcolor{rowgray}
VCI / SSI iBoard & OHLCV, khối ngoại, index & Postgres & -- \\
yfinance / CoinGecko & World, commodity, crypto, FX & Không lưu & -- \\
\bottomrule
\end{tabularx}
\end{table}

\newpage

% ============================================================
% CHƯƠNG 3: GRAPH LANGGRAPH
% ============================================================
\chapter{Graph LangGraph -- Luồng chính}

\section{LangGraph -- nền tảng điều phối}

LangGraph là framework \textbf{điều phối mức thấp}: thay vì một chuỗi gọi LLM tuyến tính, toàn bộ một lượt được mô hình thành \textbf{đồ thị trạng thái} gồm các \emph{node} (hàm Python) và \emph{edge} (cung nối). Dự án khai thác trực tiếp ba cơ chế nền tảng sau:

\begin{enumerate}
  \item \textbf{State chung -- \texttt{AgentState}}. Mọi node đọc/ghi cùng một state (\texttt{TypedDict}, \texttt{agents/state.py}). Mỗi node chỉ trả \emph{một phần} dict các trường nó thay đổi; LangGraph \textbf{merge} phần đó vào state trước khi chạy node kế tiếp. Vì state chỉ chứa text/đường dẫn (không chứa DataFrame), nó tuần tự hoá được -- điều kiện để checkpoint lưu xuống Postgres.
  \item \textbf{Điều hướng bằng hàm thuần}. Hai loại cung: \texttt{add\_edge} (cố định) và \texttt{add\_conditional\_edges} (hàm nhận state, trả \emph{key} $\rightarrow$ node đích). Mọi điểm rẽ trong dự án -- \texttt{check\_conversation}, \texttt{\_route\_after\_clarify}, \texttt{route\_after\_subqueries}, \texttt{route\_after\_critique} -- đều là \textbf{code thuần if/else}, \emph{không} gọi LLM. Đây là cốt lõi khiến chi phí mỗi lượt dự đoán được.
  \item \textbf{Checkpoint + interrupt}. \texttt{g.compile(checkpointer=PostgresCheckpointer())} khiến mỗi bước tự lưu snapshot state theo \texttt{thread\_id}. Nhờ đó \texttt{interrupt()} treo graph giữa chừng (clarify hỏi người dùng) và \texttt{Command(resume=...)} nối lại đúng chỗ; một hội thoại có thể khôi phục trạng thái sau hàng phút/giờ.
\end{enumerate}

Chọn LangGraph thay vì chuỗi prompt tuần tự vì hai lý do: \textbf{(1) luồng có vòng lặp} -- replan (quay \texttt{decompose}) và critique-retry (quay \texttt{synthesize}) là chu trình mà chuỗi tuyến tính không biểu diễn được; \textbf{(2) cần dừng/khôi phục} -- clarify xen người dùng vào giữa luồng, đòi hỏi \texttt{interrupt} + checkpoint mà chỉ graph có trạng thái bền mới đáp ứng.

\medskip\noindent\textbf{StateGraph và conditional edge -- nhìn từ \texttt{build\_graph()}}. Toàn bộ graph được dựng trong một hàm duy nhất \texttt{agents/graph.py::build\_graph()}:

{\small
\begin{verbatim}
g = StateGraph(AgentState)                   # graph chung một state
g.add_node("classify_node", classify_node)   # node = hàm Python
g.set_entry_point("classify_node")           # node bắt đầu
g.add_conditional_edges("classify_node", check_conversation,
    {"skip": END, "out_of_scope": "node_out_of_scope",
     "verify": "clarify_node"})
g.add_edge("decompose_node", "run_subqueries_node")   # cung cố định
app = g.compile(checkpointer=checkpointer)   # biên dịch + gắn checkpoint
\end{verbatim}
}

Ba chi tiết cần nắm về \texttt{add\_conditional\_edges}: nó nhận \emph{(node gốc, hàm điều kiện, bảng ánh xạ \texttt{key $\rightarrow$ node đích})}. Hàm điều kiện đọc state rồi trả \textbf{một key}; LangGraph tra bảng để chọn node kế tiếp. Bảng phải liệt kê \emph{đủ} mọi nhánh có thể -- thiếu một key sẽ bị lỗi lúc chạy. Trong dự án, hàm điều kiện luôn là \textbf{code thuần} (\texttt{if/else}), không gọi LLM: \texttt{check\_conversation} trả \texttt{skip}/\texttt{out\_of\_scope}/\texttt{verify}, \texttt{route\_after\_subqueries} trả \texttt{replan}/\texttt{synthesize}, \texttt{route\_after\_critique} trả \texttt{save}/\texttt{retry}.

\section{Tổng quan luồng}

\texttt{agents/graph.py build\_graph()} tạo \texttt{StateGraph(AgentState)}, đăng ký node, nối edge, rồi \texttt{g.compile(checkpointer=checkpointer)}. Điểm vào là \texttt{classify\_node}. Luồng đầy đủ:

\begin{figure}[H]
\centering
\tikzset{
  bignode/.style={rectangle, rounded corners=3pt, draw=#1, fill=#1!18,
    minimum height=1.15cm, text width=3.4cm, align=center, font=\small}
}
\resizebox{\textwidth}{!}{%
\begin{tikzpicture}[every node/.style={font=\small}]
  % Terminal: user input
  \node[startstop, text width=3.4cm] (start) at (0,2.3) {USER INPUT\\câu hỏi người dùng};

  % Main spine (x = 0), top to bottom
  \node[bignode=ttublue]  (cls)   at (0,0)       {classify\\intent + ticker + trace};
  \node[bignode=ttuorange] (clar) at (0,-2.3)    {clarify\\interrupt() nếu thiếu};
  \node[bignode=ttugreen] (fast)  at (-5.0,-4.6) {fast-path\\build\_single\_subtask};
  \node[bignode=ttublue]  (dec)   at (0,-4.6)    {decompose\\generate\_sub\_tasks};
  \node[bignode=ttugreen] (run)   at (0,-6.9)    {run\_subqueries\\gather\_data, no LLM};
  \node[bignode=ttugold]  (syn)   at (0,-9.2)    {synthesize\_final\\LLM stream};
  \node[bignode=ttugold]  (crit)  at (0,-11.5)   {critique\_report\\LLM self-check};
  \node[bignode=ttugreen] (fin)   at (0,-13.8)   {finalize\\turn\_end};

  % Terminal: output
  \node[startstop, text width=3.4cm] (result) at (0,-16.1) {OUTPUT\\báo cáo / trả lời (SSE)};

  % Side exits (right column)
  \node[bignode=ttured]  (oos)   at (5.2,-2.3)  {out\_of\_scope\\từ chối};
  \node[bignode=ttugreen] (brief) at (5.2,-4.6) {market\_brief\\graph con};

  % User input -> classify
  \draw[arrowt=ttugreen] (start) -- node[right, font=\tiny]{query} (cls);

  % classify: check_conversation (3 branches)
  \draw[arrowt=ttuorange] (cls) -- node[right, font=\tiny]{verify} (clar);
  \draw[arrowt=ttured] (cls.south east) -- node[above, font=\tiny, sloped]{out\_of\_scope} (oos.north);
  \draw[arrowt=ttured] (cls.west) -- node[above, font=\tiny]{skip} (-8.6,0) -- (-8.6,-16.1) -- (result.west);

  % clarify: _route_after_clarify (5 branches)
  \draw[arrowt=ttuorange] (clar.south) -- node[right, font=\tiny]{decompose} (dec.north);
  \draw[arrowt=ttublue] (clar.south west) -- node[above, font=\tiny]{simple} (fast.north);
  \draw[arrowt=ttured] (clar.east) -- node[above, font=\tiny]{out\_of\_scope} (oos.west);
  \draw[arrowt=ttured] (clar.south east) to[out=-15, in=160] node[below, font=\tiny, near start]{market\_brief} (brief.west);
  \draw[arrowt=ttured] (clar.west) -- node[above, font=\tiny]{conversation} (-7.6,-2.3) -- (-7.6,-13.8) -- (fin.west);

  % Merge to run
  \draw[arrowt=ttublue] (fast.south) |- (run.west);
  \draw[arrowt=ttuorange] (dec.south) -- (run.north);

  % run -> synthesize -> critique -> finalize
  \draw[arrowt=ttuorange] (run) -- node[right, font=\tiny]{synthesize} (syn);
  \draw[arrowt=ttuorange] (syn) -- (crit);
  \draw[arrowt=ttuorange] (crit) -- node[right, font=\tiny]{save} (fin);
  \draw[arrowt=ttuorange] (fin) -- (result);

  % Loops (dashed)
  \draw[arrowt=ttured, dashed] (run.west) to[bend left=10] node[left, font=\tiny]{replan} (dec.west);
  \draw[arrowt=ttured, dashed] (crit.east) to[bend right=10] node[right, font=\tiny]{retry} (syn.east);

  % Exits to output (right lanes)
  \draw[arrowt=ttured] (oos.east) -- (9.2,-2.3) -- (9.2,-16.1) -- (result.east);
  \draw[arrowt=ttured] (brief.east) -- (8.0,-4.6) -- (8.0,-13.8) -- (fin.east);
\end{tikzpicture}}
\caption{Luồng graph trạng thái LangGraph của trợ lý: từ câu hỏi người dùng đến báo cáo. Nhãn trên cung là \emph{key} mà hàm điều kiện trả về (\texttt{check\_conversation}, \texttt{\_route\_after\_clarify}, \texttt{route\_after\_subqueries}, \texttt{route\_after\_critique}); cung nét đứt là vòng lặp.}
\label{fig:graph}
\end{figure}

\section{Bảng node -- vai trò và LLM call}

\begin{table}[H]
\centering
\caption{Các node và nhiệm vụ}
\small
\begin{tabularx}{\textwidth}{lXl}
\toprule
\textbf{Node} & \textbf{Nhiệm vụ} & \textbf{Gọi LLM?} \\
\midrule
\texttt{classify\_node} & Gán intent + ticker + time\_context; mở trace Langfuse & Không (production) \\
\rowcolor{rowgray}
\texttt{clarify\_node} & \texttt{interrupt()} hỏi khi thiếu intent/ticker; resume thì phân loại lại & Có (phân loại lại) \\
\texttt{node\_out\_of\_scope} & Từ chối crypto/ngoại/forex bằng câu cố định & Không \\
\rowcolor{rowgray}
\texttt{node\_market\_brief} & Gọi \texttt{build\_brief\_graph()} (graph con) sinh market brief & Có (graph con) \\
\texttt{decompose\_node} & Phân rã query thành ≤4 sub-task, mỗi task có intent & Có \\
\rowcolor{rowgray}
\texttt{build\_single\_subtask\_node} & Fast-path: gói intent+ticker thành 1 sub-task, bỏ decompose & Không \\
\texttt{run\_subqueries\_node} & Chạy \texttt{gather\_data} từng sub-task; fan-out song song & \textbf{Không} \\
\rowcolor{rowgray}
\texttt{synthesize\_final} & Gộp dữ liệu thành báo cáo cuối; stream token ra SSE & Có \\
\texttt{critique\_report\_node} & LLM tự chấm báo cáo theo checklist, trả \texttt{\{pass, feedback\}} & Có \\
\rowcolor{rowgray}
\texttt{finalize\_node} & Đóng trace span của turn & Không \\
\bottomrule
\end{tabularx}
\end{table}

\section{Bảng edge -- điều hướng}

\begin{table}[H]
\centering
\caption{Conditional edges (hàm trả key → node đích)}
\small
\begin{tabularx}{\textwidth}{lX}
\toprule
\textbf{Hàm} & \textbf{Keys → đích} \\
\midrule
\texttt{check\_conversation} & \texttt{skip → END}, \texttt{out\_of\_scope → node\_out\_of\_scope}, \texttt{verify → clarify\_node} \\
\rowcolor{rowgray}
\texttt{\_route\_after\_clarify} & \texttt{market\_brief}, \texttt{simple → build\_single\_subtask}, \texttt{decompose}, \texttt{conversation}, \texttt{out\_of\_scope} \\
\texttt{route\_after\_subqueries} & \texttt{replan → decompose\_node}, \texttt{synthesize → synthesize\_final} \\
\rowcolor{rowgray}
\texttt{route\_after\_critique} & \texttt{save → finalize\_node}, \texttt{retry → synthesize\_final} \\
\texttt{\_check\_approval\_decision} (human\_approval) & \texttt{end → END}, \texttt{synthesize\_final} \\
\bottomrule
\end{tabularx}
\end{table}

\section{Điều hướng thuần logic}

Ba quyết định điều hướng cốt lõi là \textbf{code thuần} (if/else), không gọi LLM:

\begin{enumerate}
  \item \texttt{route\_after\_subqueries}: nếu tỷ lệ sub-task rỗng \texttt{> 0.5} và chưa replan → \texttt{replan} (quay \texttt{decompose\_node} đúng một lần); ngược lại \texttt{synthesize}.
  \item \texttt{route\_after\_critique}: nếu \texttt{critique\_pass} → \texttt{save}; chưa quá \texttt{MAX\_CRITIQUE} → \texttt{retry} (gấp \texttt{critique\_feedback} vào prompt synthesize); quá → \texttt{save}.
  \item \texttt{\_route\_after\_clarify}: intent phức hợp (\texttt{market\_brief}, \texttt{investment\_case}, \texttt{macro\_sector}) → \texttt{decompose}; intent leaf có ticker → \texttt{simple} (fast-path); leaf ngành (\texttt{screening}, \texttt{breakout\_scan}) → \texttt{simple}; còn lại \texttt{decompose}.
\end{enumerate}

\section{Walkthrough -- ``Đánh giá HPG hôm nay''}

\begin{table}[H]
\centering
\caption{Chuyển trạng thái theo node cho câu hỏi ``đánh giá HPG hôm nay''}
\small
\begin{tabularx}{\textwidth}{llX}
\toprule
\textbf{Bước} & \textbf{Node} & \textbf{Hành động} \\
\midrule
1 & router (ngoài graph) & LLM gọi \texttt{needs\_agent\_run(intent=investment\_case, ticker=HPG)} \\
\rowcolor{rowgray}
2 & \texttt{classify\_node} & Dùng kết quả router, mở trace \\
3 & \texttt{clarify\_node} & \texttt{detect\_ambiguity} → \texttt{None} (đủ intent+ticker) → đi tiếp \\
\rowcolor{rowgray}
4 & \texttt{\_route\_after\_clarify} & \texttt{investment\_case ∈ complex} → \texttt{decompose} \\
5 & \texttt{decompose\_node} & \texttt{generate\_sub\_tasks(n=4)} → 4 sub-task \\
\rowcolor{rowgray}
6 & \texttt{run\_subqueries\_node} & Chạy 4 \texttt{gather\_data}; intent fan-out song song \\
7 & \texttt{route\_after\_subqueries} & \texttt{empty\_ratio ≤ 0.5} → \texttt{synthesize} \\
\rowcolor{rowgray}
8 & \texttt{synthesize\_final} & 1 LLM call tổng hợp, stream ra SSE \\
9 & \texttt{critique\_report\_node} & LLM chấm checklist → \texttt{pass=True} \\
\rowcolor{rowgray}
10 & \texttt{finalize\_node} & Đóng trace \\
\bottomrule
\end{tabularx}
\end{table}

\newpage

% ============================================================
% CHƯƠNG 4: RAG-FUSION & RRF
% ============================================================
\chapter{RAG-Fusion và Reciprocal Rank Fusion}

\section{Ý tưởng}

RAG-Fusion thay truy hồi đơn lẻ bằng hai bước: (1) \textbf{sinh nhiều truy vấn con} đa góc nhìn; (2) \textbf{hợp nhất thứ hạng} từ nhiều danh sách kết quả bằng RRF. Graph độc lập \texttt{rag/rag\_fusion\_graph.py} cài đặt kỹ thuật này với 5 node:

\begin{figure}[H]
\centering
\begin{tikzpicture}[node distance=0.55cm, every node/.style={font=\small}]
  \node[process] (d) {decompose\\LLM sinh N sub-query};
  \node[process, below=0.55cm of d] (m) {multi\_retrieve\\BM25 + vector + news\\+ Postgres + Tavily};
  \node[process, below=0.55cm of m] (r) {rrf\_fuse\\RRF fold → top-K};
  \node[process, below=0.55cm of r] (a) {analyze\\LLM, guard\\không mua/bán};
  \node[startstop, below=0.55cm of a] (p) {report\\format + trích nguồn};

  \draw[arrowt=ttublue] (d) -- (m);
  \draw[arrowt=ttublue] (m) -- (r);
  \draw[arrowt=ttublue] (r) -- (a);
  \draw[arrowt=ttuorange] (a) -- (p);
\end{tikzpicture}
\caption{Graph RAG-Fusion: decompose → multi\_retrieve → rrf\_fuse → analyze → report}
\label{fig:ragfusion}
\end{figure}

\section{Sinh truy vấn con}

\texttt{rag/multi\_query.py} cung cấp hai hàm sinh truy vấn con:

\begin{keybox}{generate\_sub\_queries vs generate\_sub\_tasks}
\texttt{generate\_sub\_queries(query, n=4)}: sinh N truy vấn con thuần -- dùng cho graph RAG-Fusion.\\
\texttt{generate\_sub\_tasks(query, n=4)}: sinh N sub-task \textbf{gắn intent} (\texttt{price\_action}, \texttt{technical\_analysis}, \texttt{valuation}, \ldots) -- dùng cho graph agent chính. Cả hai đều giới hạn ≤5 mã tiêu biểu mỗi câu, lọc \texttt{<think>}, bỏ số thứ tự, và fallback về query gốc khi LLM trả rỗng.
\end{keybox}

\section{Reciprocal Rank Fusion}

\texttt{rag/fusion.py rrf\_fusion} hợp nhất hai danh sách đã chấm điểm chỉ bằng thứ hạng:

\begin{equation}
\text{score}(d) = \sum_{i} \frac{1}{k + \text{rank}_i(d)},
\qquad k = 60
\label{eq:rrf}
\end{equation}

Trong đó $\text{rank}_i(d)$ là vị trí (1-based) của tài liệu $d$ trong danh sách truy hồi thứ $i$, $k$ là hằng số làm trơn (Cormack 2009).

\begin{keybox}{Tại sao dùng thứ hạng, không cộng điểm tương đồng?}
Điểm BM25 (unbounded, phụ thuộc corpus) và cosine (bounded $0$--$1$) không cộng trực tiếp được: BM25 luôn áp đảo cosine. RRF bỏ score, chỉ dùng rank -- đại lượng cộng được xuyên danh sách, không cần chuẩn hoá, không cần huấn luyện.
\end{keybox}

RRF áp dụng ở \textbf{hai vị trí}:

\begin{enumerate}
  \item \texttt{\_hybrid\_retrieve}: BM25 + Qdrant vector của \textbf{một} sub-query → RRF → top-k (fusion trong một query).
  \item \texttt{rrf\_fuse\_node}: N danh sách kết quả của N sub-query → RRF fold tuần tự → top-k (fusion xuyên query).
\end{enumerate}

\begin{warnbox}{Ràng buộc trung lập (analyze node)}
System prompt của \texttt{analyze} chứa \texttt{GUARD 1}: \textit{``TUYỆT ĐỐI không đưa ra lời khuyên mua, bán, hoặc nắm giữ bất kỳ cổ phiếu nào.''} Kết quả web (\texttt{[WEB]}) phải gắn nhãn \textit{``Chưa được kiểm chứng''}; nguồn BCTC/tin tức phải ghi rõ nguồn.
\end{warnbox}

\newpage

% ============================================================
% CHƯƠNG 5: TẦNG TOOLS & DỮ LIỆU
% ============================================================
\chapter{Tầng tools và dữ liệu}

\section{Hợp đồng ToolResult}

Mọi tool trả \texttt{ToolResult(status, data, message)} -- \textbf{không raise, không trả list rỗng trần}. \texttt{status ∈ \{ ok, no\_data, invalid\_input, upstream\_error, rate\_limited \}}. Nhờ vậy \texttt{\_safe\_gather} chỉ cần bọc \texttt{try/except} là đủ -- lỗi tool không bao giờ làm sập node.

\section{Tầng provider giá}

\texttt{PriceProvider} (ABC) + chuỗi fallback cho cổ phiếu VN: \textbf{Fireant → VCI → TCBS}; chỉ số VN → \texttt{SsiIndexProvider}; mã ngoại → \texttt{YFinanceProvider}. \texttt{\_detect\_provider} quyết định theo hình thái ticker và universe \texttt{securities}.

\section{Intent module -- gather\_data / run}

Mỗi intent tách hai hàm:

\begin{itemize}
  \item \texttt{gather\_data}: thuần Python, trả dữ liệu text cho graph -- không gọi LLM. Đây chính là thứ làm \texttt{run\_subqueries\_node} ``không cần LLM''.
  \item \texttt{run}: có LLM, theo market-brief pattern (Python tính số → LLM viết slot → Python lắp Markdown).
\end{itemize}

\section{Dữ liệu trong graph}

Bảng \texttt{securities} (ticker, exchange, name\_aliases, tags) là universe để phân giải mã/ngành/tên. Các bảng Postgres còn lại cấp số liệu cho từng intent: \texttt{financial\_facts} (KQKD/CDKT/LCTT), \texttt{stock\_prices} + \texttt{ohlcv\_daily} (giá), \texttt{news\_articles} (tin), \texttt{market\_index\_daily} (index), \texttt{foreign\_flows} (khối ngoại), \texttt{corporate\_events} (sự kiện quyền), \texttt{broker\_views} (nhận định CTCK).

\newpage

% ============================================================
% CHƯƠNG 6: CƠ CHẾ VẬN HÀNH
% ============================================================
\chapter{Cơ chế vận hành}

\section{State -- text, không DataFrame}

\texttt{AgentState} là \texttt{TypedDict(total=False)}. Mỗi node trả một dict con (chỉ các trường nó thay đổi); LangGraph merge vào state. \textbf{State chỉ chứa text/đường dẫn, không chứa DataFrame} -- để checkpoint tuần tự hoá được và không phình.

Các trường then chốt: routing (\texttt{intent}, \texttt{ticker}, \texttt{tickers}, \texttt{sector}, \texttt{original\_query}), decompose/replan (\texttt{sub\_tasks}, \texttt{sub\_results}, \texttt{sub\_results\_empty\_ratio}, \texttt{replan\_attempted}), critique (\texttt{critique\_pass}, \texttt{critique\_feedback}, \texttt{critique\_attempts}, \texttt{report\_candidates}), ngân sách (\texttt{llm\_calls}, \texttt{turn\_started\_at}).

\section{Ngân sách turn -- chống loop vô hạn}

\begin{table}[H]
\centering
\caption{Các vòng lặp và giới hạn}
\small
\begin{tabularx}{\textwidth}{lXl}
\toprule
\textbf{Vòng lặp} & \textbf{Điều kiện} & \textbf{Giới hạn} \\
\midrule
Replan & \texttt{empty\_ratio > 0.5} và chưa replan & $\leq 1$ lần \\
\rowcolor{rowgray}
Critique retry & \texttt{critique\_pass == False} & $\leq \texttt{MAX\_CRITIQUE}$ (1) \\
Ngân sách turn & \texttt{llm\_calls ≥ 8} hoặc quá 180s & chặn mọi vòng \\
\bottomrule
\end{tabularx}
\end{table}

\texttt{\_llm\_budget\_exceeded(state)} được gọi ở mọi điểm rẽ nhánh có thể quay vòng (\texttt{route\_after\_subqueries}, \texttt{route\_after\_critique}) -- hết ngân sách thì đi thẳng đến synthesize/save, không thêm LLM call.

\section{Interrupt \& checkpoint}

\begin{itemize}
  \item \texttt{clarify\_node} gọi \texttt{interrupt(question)} treo graph khi thiếu intent/ticker; resume bằng \texttt{Command(resume=...)}.
  \item \texttt{human\_approval} (tùy chọn) chèn \texttt{request\_approval} trước synthesize để người duyệt trước khi viết báo cáo.
  \item Checkpoint: \texttt{thread\_id = conversation\_id} → mỗi hội thoại một luồng checkpoint riêng (Postgres).
\end{itemize}

\begin{warnbox}{Bug tinh tế khi resume}
Resume khôi phục state \textbf{từ thời điểm interrupt tạo} -- có thể cũ hàng phút/giờ. Hai hệ quả đã xử lý: (1) reset \texttt{turn\_started\_at}/\texttt{llm\_calls} khi resume để ngân sách turn không bị sai; (2) interrupt quá hạn (\texttt{INTERRUPT\_STALE\_SECONDS}) bị bỏ rơi bằng \texttt{thread\_id} mới.
\end{warnbox}

\section{Streaming ra SSE}

\texttt{synthesize\_final} stream token qua \texttt{emit\_llm\_delta} → observer → \texttt{asyncio.Queue} → SSE. Chỉ lượt đầu stream; critique-retry dùng \texttt{generate()} để SSE không nhận \texttt{draft1 + draft2} nối nhau.

\section{Đánh giá chất lượng -- RAGAS + Langfuse}

Hai lớp quan sát tách biệt: \textbf{Langfuse} theo dõi từng lượt chạy thật (trace/span, chi phí, độ trễ); \textbf{RAGAS} đo chất lượng câu trả lời trên nhánh \texttt{rag\_qa} bằng 4 tiêu chí chuẩn của thư viện.

\subsection{Tracing -- Langfuse}

\texttt{llm.factory.create\_client()} bọc client LLM theo thứ tự \texttt{\_wrap\_langfuse()} → \texttt{tracing.instrument\_llm()} (\texttt{llm/factory.py}). Khi \texttt{LANGFUSE\_PUBLIC\_KEY}/\texttt{LANGFUSE\_SECRET\_KEY} có trong \texttt{.env}, mọi lệnh gọi \texttt{generate()} -- kể cả \texttt{classify\_hybrid}/\texttt{llm\_classify} (phân loại intent) và \texttt{conversation\_router.llm\_route} (định tuyến) -- tự động lên trace Langfuse; thiếu key thì rơi về ghi cục bộ \texttt{traces/latest.jsonl} (\texttt{tracing.\_emit}), không lỗi. \texttt{graph\_callbacks()} (\texttt{tracing.py}) gắn thêm \texttt{langfuse.langchain.CallbackHandler} cho từng node LangGraph khi chạy \texttt{app.invoke(..., config={"callbacks": graph\_callbacks()})}.

\subsection{Chất lượng câu trả lời -- RAGAS trên nhánh rag\_qa}

Script \texttt{evals/eval\_rag\_qa\_ragas.py} đo pipeline \texttt{rag\_qa} thật, không mock: \texttt{rag.qa.retrieve\_only()} lấy ngữ cảnh (SQL rows + chunk RAG), \texttt{rag.qa.answer()} sinh câu trả lời qua DeepSeek. \texttt{ground\_truth} không hard-code -- lấy trực tiếp từ \texttt{rag.sql\_agent.execute\_safe()} trên cùng câu hỏi, tức số liệu thật trong Postgres tại thời điểm chạy.

\begin{table}[H]
\centering
\caption{4 tiêu chí RAGAS áp dụng}
\small
\begin{tabularx}{\textwidth}{lX}
\toprule
\textbf{Tiêu chí} & \textbf{Ý nghĩa} \\
\midrule
\texttt{faithfulness} & Câu trả lời có bám sát ngữ cảnh truy hồi, không bịa thêm \\
\rowcolor{rowgray}
\texttt{answer\_relevancy} & Câu trả lời có đúng trọng tâm câu hỏi \\
\texttt{context\_precision} & Ngữ cảnh truy hồi có ít nhiễu (so với \texttt{ground\_truth}) \\
\rowcolor{rowgray}
\texttt{context\_recall} & Ngữ cảnh truy hồi có đủ thông tin để trả lời đúng \\
\bottomrule
\end{tabularx}
\end{table}

\begin{warnbox}{Ràng buộc dữ liệu}
Cả 4 tiêu chí chỉ tin được khi \texttt{execute\_safe()} trả về rows thật cho câu hỏi -- script cảnh báo (\texttt{WARNING}) và in ground\_truth ra màn hình trước khi chấm điểm để soát tay, thay vì âm thầm chấm trên ground\_truth trống.
\end{warnbox}

\newpage

% ============================================================
% CHƯƠNG 7: KẾT LUẬN
% ============================================================
\chapter{Kết luận}

\section{Tổng kết Đóng góp}

\begin{table}[H]
\centering
\caption{Tóm tắt kết quả theo yêu cầu đề tài}
\small
\begin{tabularx}{\textwidth}{lXl}
\toprule
\textbf{Yêu cầu} & \textbf{Kết quả} & \textbf{Căn cứ} \\
\midrule
Graph agent & Luồng 10 node classify → finalize & Chương 3 \\
\rowcolor{rowgray}
RAG-Fusion & Graph 5 node + RRF $k=60$ & Chương 4 \\
Dữ liệu & Postgres/Qdrant/MinIO + ToolResult & Chương 5 \\
\rowcolor{rowgray}
Ràng buộc trung lập & \texttt{STRICT\_NEUTRAL} / GUARD cấm mua-bán & Chương 4 \\
Chống loop & Ngân sách 8 call / 180s, replan ≤1 & Chương 6 \\
\bottomrule
\end{tabularx}
\end{table}

\section{Ba quyết định thiết kế then chốt}

\begin{enumerate}
  \item \textbf{Điều hướng bằng code thuần} -- routing/replan/critique là if/else; chỉ node sinh nội dung gọi LLM. Chi phí và độ trễ dự đoán được.
  \item \textbf{State chỉ chứa text/đường dẫn} -- checkpoint tuần tự hoá được, và \texttt{run\_subqueries} gather dữ liệu dạng text không cần LLM.
  \item \textbf{Mọi vòng lặp có trần} -- replan ≤1, critique ≤\texttt{MAX\_CRITIQUE}, ngân sách turn 8 call/180s, đảm bảo turn ``bệnh hoạn'' không đốt tiền vô hạn.
\end{enumerate}

\section{Hướng Phát triển}

\begin{itemize}
  \item \textbf{Reranking}: thêm cross-encoder sau RRF để tinh chỉnh top-k trước khi tổng hợp.
  \item \textbf{HyDE}: sinh tài liệu giả định trước khi truy hồi để thu hẹp khoảng cách ngữ nghĩa query--corpus.
  \item \textbf{Đa nguồn đồng bộ}: hợp nhất fusion với database và web search trong cùng một node tổng hợp.
\end{itemize}

\newpage

% ============================================================
% TÀI LIỆU THAM KHẢO
% ============================================================
\begin{thebibliography}{99}

\bibitem{cormack2009}
Cormack, G.~V., Clarke, C.~L.~A., \& Buettcher, S. (2009).
Reciprocal rank fusion outperforms Condorcet and individual rank learning methods.
\textit{Proceedings of the 32nd International ACM SIGIR Conference}, 758--759.

\bibitem{langgraph}
LangChain. (2024).
\textit{LangGraph -- Low-level orchestration framework for agentic applications}.
\url{https://langchain-ai.github.io/langgraph/}

\bibitem{ragas}
Es, S., James, J., Espinosa-Anke, L., \& Schockaert, S. (2024).
RAGAS: Automated evaluation of retrieval augmented generation.
\textit{Proceedings of the 18th Conference of the European Chapter of the ACL}.

\bibitem{qdrant}
Qdrant. (2024).
\textit{Qdrant -- vector database}.
\url{https://qdrant.tech/}

\bibitem{ragfusion}
Raudaschl, A. (2023).
\textit{Forget RAG, the future is RAG-Fusion}.
\url{https://towardsdatascience.com/forget-rag-the-future-is-rag-fusion}

\bibitem{requirement}
Giảng viên Project Agent. (2026).
\textit{Hướng dẫn thực hiện Project môn Agent -- Đề tài 4: Trợ lý phân tích thị trường -- RAG-Fusion}.
Khoa Công nghệ Thông tin, Trường Đại học Tân Tạo.

\end{thebibliography}

\end{document}
